% GENERATED FILE. Edit canonical lessons, then run npm run generate:books.
% canonical-source-hash: fnv1a64:926d3f0116d9a414
% canonical-lessons: MR-C306-kramasah, MR-C306-anekda, MR-C306-adhunmadhun, MR-C306-apoap, MR-C306-nuktac

\chapter{Bit by Bit}
\label{ch:mr-bit-by-bit}
\hlchaptermodality{306}

\begin{chapteropening}
\textbf{By the end of this chapter:} \emph{I can say how often and how gradually things happened.}

Complete the last lesson of chapter 306: I can say how often and how gradually things happened.
\end{chapteropening}

\section[\texorpdfstring{\mr{क्रमशः} (kramaśaḥ)}{क्रमशः (kramaśaḥ)}]{\mr{क्रमशः} (kramaśaḥ) — step by step; gradually}

\label{lesson:MR-C306-kramasah}



\begin{quote}
Before the new one: say the Marathi for over time, then the Marathi for to elapse.
\end{quote}

\subsection*{You'll want to know: \mr{क्रमशः}}
\textbf{\mr{क्रमशः}} — \emph{kramaśaḥ} — \textquotedblleft{}step by step; gradually\textquotedblright{}.

\begin{quote}
\mr{नदीचं} \mr{पाणी} \mr{क्रमशः} \mr{कमी} \mr{झालं}.
\end{quote}

\begin{quote}
\emph{nadīcaṁ pāṇī kramaśaḥ kamī jhālaṁ.}
\end{quote}

\begin{quote}
The water in the river went down step by step.
\end{quote}

Formal; at the end of a serial story it means to be continued.

\begin{cousinweb}
\textbf{\mr{क्रमशः}} (\emph{kramaśaḥ}) is a tatsama, taken from Sanskrit \emph{kramaśas}, 'gradually, by degrees' (Monier-Williams).
\end{cousinweb}

\begin{grammarlens}[title={What you've built}]
One more word for telling what happened.
\end{grammarlens}

\subsection*{Your turn}
\begin{itemize}
\raggedright
  \item \emph{Say it:} \emph{kramaśaḥ}
  \item \emph{Say it:} \emph{kramaśaḥ}, once more
  \item \emph{Say it:} \emph{nadīcaṁ pāṇī kramaśaḥ kamī jhālaṁ.} — The water in the river went down step by step.
  \item \emph{Recall:} say the Marathi for over time, then the Marathi for to elapse, then say \emph{kramaśaḥ} again
\end{itemize}

\begin{tcolorbox}[breakable,colback=teal!4,colframe=teal!35!black,title={Before you move on}]
What does \mr{क्रमशः} mean? (\textquotedblleft{}Step by step; gradually\textquotedblright{}.) Say it once more.
\end{tcolorbox}
\section[\texorpdfstring{\mr{अनेकदा} (anekdā)}{अनेकदा (anekdā)}]{\mr{अनेकदा} (anekdā) — many times}

\label{lesson:MR-C306-anekda}



\begin{quote}
Before the new one: say the Marathi for to elapse, then the Marathi for step by step.
\end{quote}

\subsection*{You'll want to know: \mr{अनेकदा}}
\textbf{\mr{अनेकदा}} — \emph{anekdā} — \textquotedblleft{}many times\textquotedblright{}.

\begin{quote}
\mr{मी} \mr{हा} \mr{किल्ला} \mr{अनेकदा} \mr{पाहिला} \mr{आहे}.
\end{quote}

\begin{quote}
\emph{mī hā killā anekdā pāhilā āhe.}
\end{quote}

\begin{quote}
I have seen this fort many times.
\end{quote}

Built from \mr{अनेक}, many, and -\mr{दा}, times.

\begin{cousinweb}
\textbf{\mr{अनेकदा}} (\emph{anekdā}) is built from Sanskrit \emph{aneka}, 'not one, many' (Monier-Williams), and the ending \emph{-dā}, 'times'.
\end{cousinweb}

\begin{grammarlens}[title={What you've built}]
One more word for telling what happened.
\end{grammarlens}

\subsection*{Your turn}
\begin{itemize}
\raggedright
  \item \emph{Say it:} \emph{anekdā}
  \item \emph{Say it:} \emph{anekdā}, once more
  \item \emph{Say it:} \emph{mī hā killā anekdā pāhilā āhe.} — I have seen this fort many times.
  \item \emph{Recall:} say the Marathi for to elapse, then the Marathi for step by step, then say \emph{anekdā} again
\end{itemize}

\begin{tcolorbox}[breakable,colback=teal!4,colframe=teal!35!black,title={Before you move on}]
What does \mr{अनेकदा} mean? (\textquotedblleft{}Many times\textquotedblright{}.) Say it once more.
\end{tcolorbox}
\section[\texorpdfstring{\mr{अधूनमधून} (adhūnmadhūn)}{अधूनमधून (adhūnmadhūn)}]{\mr{अधूनमधून} (adhūnmadhūn) — now and then}

\label{lesson:MR-C306-adhunmadhun}



\begin{quote}
Before the new one: say the Marathi for step by step, then the Marathi for many times.
\end{quote}

\subsection*{You'll want to know: \mr{अधूनमधून}}
\textbf{\mr{अधूनमधून}} — \emph{adhūnmadhūn} — \textquotedblleft{}now and then\textquotedblright{}.

\begin{quote}
\mr{आजोबा} \mr{अधूनमधून} \mr{जुन्या} \mr{गोष्टी} \mr{सांगतात}.
\end{quote}

\begin{quote}
\emph{ājobā adhūnmadhūn junyā goṣṭī sāṅgtāt.}
\end{quote}

\begin{quote}
Grandpa tells old stories now and then.
\end{quote}

A rhyming pair; \mr{कधीकधी} means much the same.

\begin{cousinweb}
\textbf{\mr{अधूनमधून}} (\emph{adhūnmadhūn}) is a rhyming pair built on \emph{madhūn}, 'from the middle'; the origin of the first half, \emph{adhūn}, is not settled.
\end{cousinweb}

\begin{grammarlens}[title={What you've built}]
One more word for telling what happened.
\end{grammarlens}

\subsection*{Your turn}
\begin{itemize}
\raggedright
  \item \emph{Say it:} \emph{adhūnmadhūn}
  \item \emph{Say it:} \emph{adhūnmadhūn}, once more
  \item \emph{Say it:} \emph{ājobā adhūnmadhūn junyā goṣṭī sāṅgtāt.} — Grandpa tells old stories now and then.
  \item \emph{Recall:} say the Marathi for step by step, then the Marathi for many times, then say \emph{adhūnmadhūn} again
\end{itemize}

\begin{tcolorbox}[breakable,colback=teal!4,colframe=teal!35!black,title={Before you move on}]
What does \mr{अधूनमधून} mean? (\textquotedblleft{}Now and then\textquotedblright{}.) Say it once more.
\end{tcolorbox}
\section[\texorpdfstring{\mr{आपोआप} (āpoāp)}{आपोआप (āpoāp)}]{\mr{आपोआप} (āpoāp) — by itself; automatically}

\label{lesson:MR-C306-apoap}



\begin{quote}
Before the new one: say the Marathi for many times, then the Marathi for now and then.
\end{quote}

\subsection*{You'll want to know: \mr{आपोआप}}
\textbf{\mr{आपोआप}} — \emph{āpoāp} — \textquotedblleft{}by itself; automatically\textquotedblright{}.

\begin{quote}
\mr{वारा} \mr{आला} \mr{आणि} \mr{दार} \mr{आपोआप} \mr{बंद} \mr{झालं}.
\end{quote}

\begin{quote}
\emph{vārā ālā āṇi dār āpoāp band jhālaṁ.}
\end{quote}

\begin{quote}
The wind blew and the door shut by itself.
\end{quote}

A doubling of \mr{आप}, self.

\begin{cousinweb}
\textbf{\mr{आपोआप}} (\emph{āpoāp}) is built from \emph{āp}, 'self', said twice; CDIAL lists Marathi \emph{āplā}, 'own', under Sanskrit \emph{ātman}, 'self' (CDIAL 1135).
\end{cousinweb}

\begin{grammarlens}[title={What you've built}]
One more word for telling what happened.
\end{grammarlens}

\subsection*{Your turn}
\begin{itemize}
\raggedright
  \item \emph{Say it:} \emph{āpoāp}
  \item \emph{Say it:} \emph{āpoāp}, once more
  \item \emph{Say it:} \emph{vārā ālā āṇi dār āpoāp band jhālaṁ.} — The wind blew and the door shut by itself.
  \item \emph{Recall:} say the Marathi for many times, then the Marathi for now and then, then say \emph{āpoāp} again
\end{itemize}

\begin{tcolorbox}[breakable,colback=teal!4,colframe=teal!35!black,title={Before you move on}]
What does \mr{आपोआप} mean? (\textquotedblleft{}By itself; automatically\textquotedblright{}.) Say it once more.
\end{tcolorbox}
\section[\texorpdfstring{\mr{नुकताच} (nuktāc)}{नुकताच (nuktāc)}]{\mr{नुकताच} (nuktāc) — lately; a moment ago}

\label{lesson:MR-C306-nuktac}



\begin{quote}
Before the new one: say the Marathi for now and then, then the Marathi for by itself.
\end{quote}

\subsection*{You'll want to know: \mr{नुकताच}}
\textbf{\mr{नुकताच}} — \emph{nuktāc} — \textquotedblleft{}lately; a moment ago\textquotedblright{}.

\begin{quote}
\mr{तो} \mr{नुकताच} \mr{गावाहून} \mr{परत} \mr{आला}.
\end{quote}

\begin{quote}
\emph{to nuktāc gāvāhūn parat ālā.}
\end{quote}

\begin{quote}
He has recently come back from the village.
\end{quote}

It agrees like an adjective: \mr{तो} \mr{नुकताच} \mr{आला}, \mr{ती} \mr{नुकतीच} \mr{आली}, \mr{ते} \mr{नुकतंच} \mr{आलं}.

\begin{cousinweb}
\textbf{\mr{नुकताच}} (\emph{nuktāc}) is a Marathi word for lately, made with the emphatic \emph{-c}; its origin is not settled.
\end{cousinweb}

\begin{grammarlens}[title={What you've built}]
One more word for telling what happened.
\end{grammarlens}

\subsection*{Your turn}
\begin{itemize}
\raggedright
  \item \emph{Say it:} \emph{nuktāc}
  \item \emph{Say it:} \emph{nuktāc}, once more
  \item \emph{Say it:} \emph{to nuktāc gāvāhūn parat ālā.} — He has recently come back from the village.
  \item \emph{Recall:} say the Marathi for now and then, then the Marathi for by itself, then say \emph{nuktāc} again
\end{itemize}

\begin{tcolorbox}[breakable,colback=teal!4,colframe=teal!35!black,title={Before you move on}]
What does \mr{नुकताच} mean? (\textquotedblleft{}Lately; a moment ago\textquotedblright{}.) Say it once more.
\end{tcolorbox}
